% Generated TeX body; no PDF or document-style claim.
% Maintained source: problems/degree-constrained-subgraphs/reference/correctness.md
% Recovered-source SHA-256: 860d6cd071022f3de2558ebcd9581bc4d11b2e9e1d934bccf454bc2676816cbd
% Rebuild: python3 tools/build_manuscripts.py
% Common packages: amsmath, amssymb, hyperref, longtable, booktabs, array.
% Shared original Subset Sum preamble: manuscripts/paper-style.tex.

\hypertarget{af-5f2ce72611-correctness-and-certificates-for-the-executable-degree-solver}{%
\section{Correctness and certificates for the executable degree solver}\label{af-5f2ce72611-correctness-and-certificates-for-the-executable-degree-solver}}

\textbf{Ready for review: unconditional local construction and ordinary polynomial-time implementation.} The code uses the complete sparse reduction in \href{https://github.com/icekylinx/algorithm-frontier/blob/main/problems/degree-constrained-subgraphs/proof.md}{Theorem 1.1 and Sections 2--5}, followed by the repository's \href{https://github.com/icekylinx/algorithm-frontier/blob/main/problems/general-matching/reference/README.md}{exact matching solver}. No numerical cycle-search or randomized primitive is used. This note identifies the implemented decisions, the exact certificate meaning and the runtime charged to them.

\hypertarget{af-5f2ce72611-native-records-and-normalization}{%
\subsection{1. Native records and normalization}\label{af-5f2ce72611-native-records-and-normalization}}

Every input edge has two distinct vertices and a stable input index. Parallel records receive separate incidences and remain independently selectable. Loops are rejected because the stated degree reduction is for loopless graphs; a loop must not silently disappear from a degree request. Supplied IDs are distinct integers or strings and remain passive labels during solving.

For every vertex, scan its incident degree \texttt{d(v)}, keep the lower bound \texttt{g(v)} unchanged and replace its upper bound by \texttt{f(v)=min(f(v),d(v))}. If \texttt{g(v)\textgreater{}f(v)}, the request is infeasible. Intersect the requested cardinality interval with \texttt{{[}0,m{]}}; an empty intersection is infeasible. All these conclusions are directly rechecked by the verifier.

After this scan, discard degree-zero vertices from every nontrivial reduction. A discarded vertex has lower bound zero, or the request has already been rejected. The active vertex count is at most \texttt{2m}. The original vertex universe is retained only for the native input and final degree vector. Optimization never rescans all isolates on each matching query.

\hypertarget{af-5f2ce72611-from-a-cardinality-interval-to-a-prescribed-factor}{%
\subsection{2. From a cardinality interval to a prescribed factor}\label{af-5f2ce72611-from-a-cardinality-interval-to-a-prescribed-factor}}

For a query interval \texttt{{[}L,U{]}}, the first hub is exactly the one in Section 5 of the main proof. Put \texttt{F=sum\_v\ f(v)} and \texttt{D=sum\_v(f(v)-g(v))}. Add \texttt{f(v)-g(v)} distinct edges from vertex \texttt{v} to a new hub, prescribe degree \texttt{f(v)} at every old vertex, and give the hub the degree interval \texttt{{[}max(0,F-2U),\ min(D,F-2L){]}}.

The key identity is \texttt{degree(hub)=F-2\textbar{}X\textbar{}}, where \texttt{X} is the selected original subset. Thus the new hub interval is equivalent to the requested cardinality condition. Its emptiness is another exact local NO. For the unconstrained interval \texttt{{[}0,m{]}}, this first augmentation is omitted.

Convert the remaining interval-degree instance to a prescribed-degree instance by Section 4: add one second hub, the required dummy incidences, and private triangles whose contributions to the hub are zero or two. Choose the hub target to have the same parity as the prescribed degree sum. This realizes every permitted total deficiency and only those deficiencies. When all degree intervals are already singletons, this second augmentation is omitted. An odd prescribed degree sum, or a prescribed degree larger than the available incident degree in the augmented graph, is an immediate NO.

All added edges are explicitly represented. Their count is linear in the original \texttt{m}, because every retained original upper bound is at most its incident degree. The construction never expands an arbitrarily large numeric input bound into that many records.

\hypertarget{af-5f2ce72611-sparse-selectors-and-native-decoding}{%
\subsection{3. Sparse selectors and native decoding}\label{af-5f2ce72611-sparse-selectors-and-native-decoding}}

For each augmented edge create two incidence ports and one cross edge between them. For an augmented vertex of degree \texttt{d} and prescribed degree \texttt{f}, exactly \texttt{d-f} ports must be matched locally.

The code builds the recursive permutation network proved in Section 3 of the main proof, splits every network position into an in/out pair, and attaches exactly \texttt{d-f} output sinks. The count of matched in and out endpoints cancels every internal transition. In a perfect matching, exactly the number of ports corresponding to those sinks uses local attachments. Therefore exactly \texttt{f} incidence ports use cross edges. Conversely, a prescribed factor determines a partial input-output permutation; the network routes it, and unused positions use idle edges. This proves both directions for the actual constructed graph.

One implemented simplification omits a selector when \texttt{d-f=0}. The same count identity forces all its ports to use cross edges. Its unused network positions could all use idle edges, so deleting that selector preserves both extension and decoding. This also handles degree-zero vertices without constructing a network of width zero.

Distinct native records always have distinct port pairs. The reduced graph is therefore simple even when the input has parallel edges, and every reduced vertex has degree at most three. Network input and output positions have no incoming and outgoing transitions, respectively, so their external attachments respect that degree bound.

Each cross edge carries either its original native index or a dummy marker. Decoding retains only matched cross edges with a native index; all hub, triangle, selector and other auxiliary edges are discarded. The original ID and endpoint orientation are then restored from that index. An independent native scan checks all vertex degrees and the selected cardinality before any \texttt{YES} leaves the solver.

\hypertarget{af-5f2ce72611-exact-feasibility-and-independently-checked-no-results}{%
\subsection{4. Exact feasibility and independently checked NO results}\label{af-5f2ce72611-exact-feasibility-and-independently-checked-no-results}}

Let the constructed matching graph have \texttt{N} vertices. The exact matcher returns a matching and its independent optimality certificate. The solver checks that certificate even when it is about to return a NO.

If the checked cardinality is \texttt{N/2}, the matching is perfect and the preceding equivalence supplies a feasible native subset. If twice its cardinality is smaller than \texttt{N}, the certificate proves that no perfect matching exists. The equivalence then proves infeasibility of the native request. An unsuccessful matching attempt without an optimality certificate would not establish this conclusion.

The independent matching verifier uses a native matching lower bound and a tight odd-component upper bound, as proved in the \href{https://github.com/icekylinx/algorithm-frontier/blob/main/problems/general-matching/reference/certificate-proof.md}{matching certificate note}. The degree verifier reconstructs the same reduction and invokes that checker. It imports no matching search. It shares the deterministic reduction builder and normalization code, so its independence is from search, not from every line of construction code.

Two easy successful cases bypass the matching call: an empty native subset when every lower degree is zero and the interval contains zero; and the whole native graph when every upper degree permits it and the interval contains \texttt{m}. Both witnesses still receive the independent native degree/cardinality check.

\hypertarget{af-5f2ce72611-optimization-and-one-interval-certificates}{%
\subsection{5. Optimization and one-interval certificates}\label{af-5f2ce72611-optimization-and-one-interval-certificates}}

First solve the full permitted interval \texttt{{[}K\_min,K\_max{]}}. If that query is infeasible, report its checked NO.

For maximization, the predicate ``some feasible subset has size in \texttt{{[}K,K\_max{]}}'' is monotone in \texttt{K}. Binary search retains one successful native witness and one final failed interval. At termination the best size \texttt{k} is supported by the witness. Unless \texttt{k=K\_max}, the retained failed query is exactly \texttt{{[}k+1,K\_max{]}}; its certificate excludes every better permitted cardinality.

For minimization, use the monotone predicate for \texttt{{[}K\_min,K{]}}. The final witness has size \texttt{k}; unless \texttt{k=K\_min}, the failed interval is exactly \texttt{{[}K\_min,k-1{]}}. The production code explicitly rejects any disagreement between the returned native size and the binary-search boundary.

The external verifier checks just the final native witness and that one excluded interval, or the relevant endpoint equality. It need not replay the binary search. No assertion about intervals of feasible \textbf{exact} cardinalities is used.

\hypertarget{af-5f2ce72611-ordinary-polynomial-cost}{%
\subsection{6. Ordinary polynomial cost}\label{af-5f2ce72611-ordinary-polynomial-cost}}

Set \texttt{H=1+m\ log(m+2)}. Both hub augmentations use \texttt{O(m+1)} records. A power-of-two network of width \texttt{s} has \texttt{s(2\ log\_2(s)+1)} positions; each chosen width is less than twice its positive incident degree. Summing over all augmented vertices gives \texttt{O(H)} matching vertices and edges, constructible in \texttt{O(H)} work after the initial native scan.

The existing exact matcher, including its certificate extraction, costs at most \texttt{O(H\^{}3)} on this graph. Its comparison sorting and independent certificate check fit that bound. Therefore a feasibility request costs \texttt{O(n+H\^{}3)}. Optimization makes at most an initial query and a logarithmic number of binary-search queries, giving \texttt{O(n+H\^{}3\ log(m+2))}. Failed interval certificates and current witnesses occupy \texttt{O(H)} space; the full native input/output occupies \texttt{O(n+m)}. Total storage is \texttt{O(n+H)}.

The external verifier reconstructs at most one reduced graph and uses comparison sorting plus graph scans, so it costs \texttt{O(n+H\ log(H+2))} time and \texttt{O(n+H)} space. It does not call the cubic matcher. Actual reading, numeric arithmetic and label-comparison costs must also be charged for representations longer than input words.

These bounds are ordinary polynomial guarantees of the shipped code. They do not discharge or reuse the faster theorem's unimplemented cycle-search requirements.

\hypertarget{af-5f2ce72611-finite-implementation-evidence}{%
\subsection{7. Finite implementation evidence}\label{af-5f2ce72611-finite-implementation-evidence}}

\href{https://github.com/icekylinx/algorithm-frontier/blob/main/problems/degree-constrained-subgraphs/verify/check_reference.py}{check\_reference.py} enumerates original edge subsets independently of the reduction and matcher. Its default suite checks 8,292 feasibility/optimization requests, including every valid degree interval on all simple graphs through three vertices and two parallel-record graphs, all cardinality ranges there, all simple four-vertex graphs with additional degree requests, and 120 seeded native cases. It invokes the matcher 5,341 times and checks 3,364 actual reduced graphs for simplicity and maximum degree three.

The suite additionally rejects 15 malformed requests and 15 corrupted certificates, checks an isolated vertex universe and long numeric/ID representations, and exercises both production command-line programs with Python assertions disabled. Production acceptance conditions are explicit checks; the mathematical test script itself refuses optimized execution. These finite results support the implementation review and do not replace the proofs above.

\textbf{Attribution:} icekylinx · GPT-6 Astra Pro.
